
This section specifies the governance structure, upgrade capabilities, and admin powers. Most protocol-controlled contracts are upgradeable; PT, YT, Market, RouterStatic, and PyLpOracle are not. Users must evaluate admin trust before depositing.

\paragraph{Authority model.}
All upgradeable contracts use OpenZeppelin's \texttt{OwnableComponent} and \texttt{UpgradeableComponent}. A single authority (the \emph{owner}) holds \texttt{DEFAULT\_ADMIN\_ROLE} and can:
\begin{itemize}
\item Replace any contract's class hash via \texttt{upgrade(new\_class\_hash)}.
\item Grant and revoke roles to other addresses.
\item Execute contract-specific admin functions.
\end{itemize}

\noindent\textbf{Target configuration:} 3-of-5 multisig with 48-hour timelock (not yet deployed).

\paragraph{Upgrade semantics.}
Starknet contracts are upgraded by replacing the \emph{class hash}---the on-chain identifier of the contract's bytecode. When \texttt{upgrade(new\_class\_hash)} executes:
\begin{enumerate}
\item The contract's class hash is atomically replaced.
\item Storage is preserved; the new implementation must be storage-compatible.
\item An \texttt{Upgraded(new\_class\_hash)} event is emitted.
\item All subsequent calls execute the new implementation immediately.
\end{enumerate}
There is no proxy indirection; the contract address remains constant while its code changes.


\section{Contract Classification}

Horizon distinguishes between \emph{internal} contracts (protocol-controlled) and \emph{external dependencies} (third-party infrastructure).

\subsection{Internal Contracts (Protocol-Controlled)}

\begin{center}
\begin{tabularx}{\textwidth}{l l >{raggedright\arraybackslash}X l}
\toprule
\textbf{Contract} & \textbf{Upgradeable} & \textbf{Scope of Upgrade Impact} & \textbf{Risk} \\
\midrule
SY & Yes & Share accounting, deposit/redeem logic & Critical \\
SYWithRewards & Yes & SY accounting plus reward distribution & Critical \\
PT & No & Redemption logic, mint authority validation & Critical \\
YT & No & Yield distribution, interest index accounting & Critical \\
Market & No & AMM curve, swap execution, fee calculation & Critical \\
Factory & Yes & Templates for new SY/PT/YT deployments & High \\
MarketFactory & Yes & Templates for new Market deployments & High \\
Router & Yes & Routing logic, could insert hidden fees & High \\
RouterStatic & No & Read-only quotes and previews & Low \\
PyLpOracle & No & PT/YT/LP TWAP pricing helper & Medium \\
PragmaIndexOracle & Yes & Index source configuration, staleness thresholds & High \\
\bottomrule
\end{tabularx}
\end{center}

\paragraph{Risk interpretation.}
\begin{description}
\item[Critical:] Upgrade can directly drain or freeze user funds.
\item[High:] Upgrade can extract value or disrupt operations but cannot drain existing positions.
\end{description}

\subsection{External Dependencies (Third-Party)}

\begin{center}
\begin{tabular}{l l l}
\toprule
\textbf{Dependency} & \textbf{Trust Requirement} & \textbf{Failure Mode} \\
\midrule
Starknet sequencer & Liveness, censorship resistance & Transactions not processed \\
Pragma oracle & Data integrity, freshness & Incorrect yield accounting \\
Underlying protocols (e.g., Nostra) & Solvency, correct yield accrual & SY value impairment \\
\bottomrule
\end{tabular}
\end{center}

\paragraph{Dependency scope.}
The protocol cannot upgrade or control external dependencies. Users accept these as infrastructure-level trust assumptions. Failures in external dependencies can cause losses even if protocol code is correct.


\section{Admin Powers}

Most admin functions require \texttt{DEFAULT\_ADMIN\_ROLE} unless otherwise noted (e.g., PragmaIndexOracle uses OPERATOR\_ROLE and PAUSER\_ROLE for specific controls). This section enumerates the exact on-chain powers.

\subsection{Powers Available on All Upgradeable Contracts}

Every contract with \texttt{OwnableComponent} and \texttt{UpgradeableComponent} exposes:

\begin{center}
\begin{tabularx}{\textwidth}{l >{raggedright\arraybackslash}X l}
\toprule
\textbf{Function} & \textbf{Effect} & \textbf{Risk} \\
\midrule
\texttt{upgrade(class\_hash)} & Replace contract implementation atomically & Critical \\
\texttt{transfer\_ownership(new\_owner)} & Transfer admin control to new address & Critical \\
\texttt{renounce\_ownership()} & Permanently remove admin control (irreversible) & N/A \\
\bottomrule
\end{tabularx}
\end{center}

\subsection{Contract-Specific Admin Functions}

\paragraph{Factory.}
\begin{center}
\begin{tabularx}{\textwidth}{l >{raggedright\arraybackslash}X l}
\toprule
\textbf{Function} & \textbf{Effect} & \textbf{Risk} \\
\midrule
\texttt{set\_sy\_class\_hash} & Change template for future SY deployments & High \\
\texttt{set\_pt\_class\_hash} & Change template for future PT deployments & Critical \\
\texttt{set\_yt\_class\_hash} & Change template for future YT deployments & Critical \\
\bottomrule
\end{tabularx}
\end{center}

\paragraph{Note.} Class hash changes affect only \emph{new} deployments. Existing SY/PT/YT instances retain their original implementation until individually upgraded.

\paragraph{MarketFactory.}
\begin{center}
\begin{tabularx}{\textwidth}{l >{raggedright\arraybackslash}X l}
\toprule
\textbf{Function} & \textbf{Effect} & \textbf{Risk} \\
\midrule
\texttt{set\_market\_class\_hash} & Change template for future Market deployments & Critical \\
\texttt{set\_treasury} & Update reserve fee recipient & High \\
\texttt{set\_default\_reserve\_fee\_percent} & Configure default reserve fee percentage & Medium \\
\texttt{set\_override\_fee} & Set per-router per-market fee override & Medium \\
\texttt{set\_default\_rate\_impact\_sensitivity} & Configure rate impact fee sensitivity & Medium \\
\texttt{set\_yield\_contract\_factory} & Update PT validation source & Medium \\
\bottomrule
\end{tabularx}
\end{center}

\paragraph{Router.}
\begin{center}
\begin{tabularx}{\textwidth}{l >{raggedright\arraybackslash}X l}
\toprule
\textbf{Function} & \textbf{Effect} & \textbf{Risk} \\
\midrule
\texttt{pause()} & Block all router operations & Medium \\
\texttt{unpause()} & Resume router operations & Low \\
\bottomrule
\end{tabularx}
\end{center}

\paragraph{PragmaIndexOracle.}
\begin{center}
\begin{tabularx}{\textwidth}{l >{raggedright\arraybackslash}X l}
\toprule
\textbf{Function} & \textbf{Effect} & \textbf{Risk} \\
\midrule
\texttt{set\_config} & Update TWAP window and staleness limits & Medium \\
\texttt{pause} / \texttt{unpause} & Freeze/unfreeze oracle reads & Medium \\
\texttt{emergency\_set\_index} & Override stored index (recovery) & High \\
\texttt{initialize\_rbac} & One-time RBAC initialization after upgrade & Medium \\
\bottomrule
\end{tabularx}
\end{center}

\paragraph{Token Contracts (SY, SYWithRewards).}
Token contracts expose only \texttt{upgrade}. Minting/burning authority is enforced by code (YT contract controls PT/YT minting), not by admin-configurable roles.

\paragraph{PT, YT, and Market.}
These contracts are not upgradeable. Their parameters are set at deployment and cannot be modified post-deployment.


\section{Role-Based Access Control}

The protocol uses OpenZeppelin's \texttt{AccessControlComponent} with three roles:

\begin{center}
\begin{tabularx}{\textwidth}{l >{raggedright\arraybackslash}X l}
\toprule
\textbf{Role} & \textbf{Permissions} & \textbf{Holder} \\
\midrule
\texttt{DEFAULT\_ADMIN\_ROLE} & All admin functions; grant/revoke other roles & Owner (multisig target) \\
\texttt{PAUSER\_ROLE} & \texttt{pause()}, \texttt{unpause()} on Router and PragmaIndexOracle & Owner (initially) \\
\texttt{OPERATOR\_ROLE} & PragmaIndexOracle configuration updates & Owner (initially) \\
\bottomrule
\end{tabularx}
\end{center}

\paragraph{Role hierarchy.}
\texttt{DEFAULT\_ADMIN\_ROLE} is the admin of all roles, including itself. Only addresses with \texttt{DEFAULT\_ADMIN\_ROLE} can grant or revoke any role.


\section{Emergency Procedures}

\subsection{Emergency Powers}

The following actions are available during emergencies:

\begin{center}
\begin{tabularx}{\textwidth}{l >{raggedright\arraybackslash}X}
\toprule
\textbf{Action} & \textbf{Effect} \\
\midrule
Pause Router & Blocks new swaps, mints, and complex operations via Router \\
Pause PragmaIndexOracle & Freezes index updates at the stored watermark \\
Emergency upgrade & Replace contract implementation without timelock (if active exploit) \\
\bottomrule
\end{tabularx}
\end{center}

\subsection{Pause Scope}

Pause mechanisms exist on the Router and PragmaIndexOracle. Router pause stops complex protocol interactions, while oracle pause freezes index updates to the stored watermark.

\begin{center}
\begin{tabularx}{\textwidth}{l l >{raggedright\arraybackslash}X}
\toprule
\textbf{Operation} & \textbf{During Pause} & \textbf{Rationale} \\
\midrule
Router swaps & Blocked & Prevent exploitation via complex operations \\
Router minting (PT/YT) & Blocked & Prevent new positions during incident \\
Router liquidity add/remove & Blocked & Prevent position changes \\
\midrule
Direct SY deposit/redeem & Available & Users can exit to underlying \\
Direct PT transfer & Available & Users can move positions \\
Direct YT transfer & Available & Users can move positions \\
Post-expiry PT redemption & Available & Users can redeem matured positions \\
YT interest claim & Available & Users can claim accrued yield \\
Market direct swap & Available & Direct contract calls bypass Router \\
LP burn (withdraw liquidity) & Available & Users can exit LP positions \\
\bottomrule
\end{tabularx}
\end{center}

\begin{center}
\fbox{\parbox{0.9\textwidth}{
\textbf{User Safety Invariant During Emergencies}

Even when the Router is paused, users can always:
\begin{enumerate}
\item Redeem SY for underlying assets (direct call to SY contract).
\item Redeem PT for SY after expiry (direct call to YT contract).
\item Claim accrued YT interest (direct call to YT contract).
\item Withdraw LP positions by burning LP tokens (direct call to Market).
\item Transfer all token types to other addresses.
\end{enumerate}
No emergency action can prevent users from exiting existing positions.
}}
\end{center}

\subsection{Emergency Response Process}

\begin{enumerate}
\item \textbf{Detection.} Exploit identified via monitoring, user report, or external disclosure.
\item \textbf{Triage.} Assess severity and funds at risk. Classify as Critical/High/Medium/Low.
\item \textbf{Containment.} If Critical or High: immediately pause Router.
\item \textbf{Remediation.} Develop and test fix. For Critical exploits, emergency upgrade without timelock.
\item \textbf{Disclosure.} Publish incident report within 24 hours of resolution.
\item \textbf{Post-mortem.} Conduct root cause analysis; implement preventive measures.
\end{enumerate}

\paragraph{Timelock bypass justification.}
Emergency upgrades (bypassing any future timelock) are permitted only when:
\begin{itemize}
\item Active exploitation is occurring or imminent.
\item Funds at risk exceed a material threshold.
\item Waiting for timelock would result in greater user losses.
\end{itemize}


\section{Upgrade Process}

\subsection{Standard Upgrade Steps}

\begin{enumerate}
\item \textbf{Development.} Implement changes; ensure storage compatibility.
\item \textbf{Testing.} Full test suite passes; invariant fuzzing on modified paths.
\item \textbf{Audit.} External audit for significant changes (new features, security-relevant modifications).
\item \textbf{Proposal.} Publish upgrade proposal with:
\begin{itemize}
\item New class hash
\item Code diff or audit report link
\item Rationale and risk assessment
\end{itemize}
\item \textbf{Timelock.} (When implemented) Submit upgrade transaction; wait for timelock period.
\item \textbf{Execution.} Multisig signers approve and execute upgrade.
\item \textbf{Verification.} Confirm new class hash matches proposal; verify contract behavior.
\end{enumerate}

\subsection{Storage Compatibility Requirements}

Upgrades must preserve storage layout:
\begin{itemize}
\item Existing storage variables retain their slots.
\item New variables are appended (not inserted).
\item Variable types are not changed.
\item Mappings and structs maintain compatible layouts.
\end{itemize}

\paragraph{Incompatible upgrades.}
If storage layout must change, a migration pattern is required: deploy new contract, migrate state, update references. This is a breaking change requiring extended user notice.


\section{Governance Attack Surface}

\begin{center}
\begin{tabularx}{\textwidth}{l l l >{raggedright\arraybackslash}X}
\toprule
\textbf{Attack Vector} & \textbf{Likelihood} & \textbf{Impact} & \textbf{Mitigations} \\
\midrule
Admin key compromise & Low & Critical (total loss) & Hardware wallet; multisig (planned); timelock (planned) \\
Malicious insider upgrade & Low & Critical (total loss) & Multisig threshold; code review; audit requirement \\
Social engineering of signers & Medium & Critical & Dedicated signing devices; verification procedures \\
Compromised upgrade proposal & Low & Critical & Class hash verification; multiple reviewer sign-off \\
Role escalation & Very Low & High & RBAC structure; role change events; monitoring \\
Pause abuse (griefing) & Low & Medium (temporary) & PAUSER\_ROLE separation; unpause requires same authority \\
\bottomrule
\end{tabularx}
\end{center}

\paragraph{Attack surface reduction.}
The governance attack surface is proportional to:
\begin{itemize}
\item Number of addresses with admin roles (currently 1; target 5 multisig signers).
\item Frequency of upgrades (minimize by shipping stable code).
\item Complexity of upgrade process (reduce via automation and verification tooling).
\end{itemize}


\section{Trust Assumptions}

Users must evaluate what they trust before using the protocol.

\subsection{Trusted (Cannot Be Verified)}

\begin{itemize}
\item Admin keys are not compromised.
\item Multisig signers do not collude maliciously.
\item Upgrades are deployed with benevolent intent.
\item Emergency actions are taken in good faith.
\item Operational security of key holders is maintained.
\end{itemize}

\subsection{Verifiable On-Chain}

\begin{itemize}
\item All contract code (class hashes can be verified against source).
\item All role grants and revocations (events emitted).
\item All upgrades (events emitted with new class hash).
\item Current pause state.
\item All token balances and positions.
\item AMM state and pool reserves.
\end{itemize}

\subsection{Trustless (Given Current Code)}

\begin{itemize}
\item Mathematical correctness of AMM pricing and yield accounting.
\item Token conservation (no inflation without corresponding deposits).
\item Access control enforcement (only authorized addresses can call admin functions).
\item Expiry mechanics (PT redeems 1:1 after expiry, YT stops accruing).
\end{itemize}

\paragraph{Key insight.}
Upgradeability means ``trustless'' properties hold only for the \emph{current} implementation. A malicious upgrade can violate any property. Users must trust admin or verify each upgrade.


\section{Governance Roadmap}

The following governance features are planned but not yet implemented:

\begin{description}
\item[Multisig (3-of-5).] Require multiple independent signers for all admin actions. Prevents single-key compromise from causing total loss.

\item[Timelock (48 hours).] Delay between upgrade proposal and execution. Provides users time to review changes and exit if concerned.

\item[Upgrade monitoring.] Automated alerts when upgrade transactions are submitted. Enables community review during timelock window.

\item[Verified class hash registry.] On-chain registry of audited, approved class hashes. Upgrades restricted to pre-verified implementations.
\end{description}

\paragraph{Decentralization considerations.}
Full decentralization (governance token, on-chain voting) is not planned for initial deployment. Rationale:
\begin{itemize}
\item Governance tokens introduce new attack vectors (vote buying, governance attacks).
\item Early-stage protocols benefit from rapid iteration capability.
\item Decentralization is meaningful only with sufficient distribution and participation.
\end{itemize}

\noindent The protocol may consider decentralized governance once the codebase is stable and user base is established.
